\documentclass[10pt]{scrartcl}
\usepackage[mdthm,colorsec,noauthor,noasy]{evan}
\usepackage{handout}
\newcommand{\vv}[1]{\mathbf{#1}}
\DeclareMathOperator{\spn}{Span}\DeclareMathOperator{\Nul}{Nul}\DeclareMathOperator{\Col}{Col}\DeclareMathOperator{\Row}{Row}
\title{Math 54 --- Week 6}
\subtitle{Solutions}
\author{}\date{}\ihead{}
\begin{document}
\maketitle

\begin{prob}\textbf{P1.}\par\idea{One row reduction gives all three bases.}
\begin{steps}
  \ii $R_2+R_1$: $(0,-2,-5,-3)$. $R_3-\tfrac32R_1$: $(0,2,5,3)$. $R_3+R_2$: zero row. Echelon form $\begin{bmatrix}-2&4&-2&-4\\0&-2&-5&-3\\0&0&0&0\end{bmatrix}$; reduced: $\begin{bmatrix}1&0&6&5\\0&1&5/2&3/2\\0&0&0&0\end{bmatrix}$.
  \ii $\Col A$: pivot columns 1, 2 of $A$: $(-2,2,-3)$, $(4,-6,8)$.
  \ii $\Row A$: $(-2,4,-2,-4)$, $(0,-2,-5,-3)$.
  \ii $\Nul A$: $x_1=-6x_3-5x_4$, $x_2=-\tfrac52x_3-\tfrac32x_4$: basis $(-6,-\tfrac52,1,0)$, $(-5,-\tfrac32,0,1)$ (or the multiples $(-12,-5,2,0)$, $(-10,-3,0,2)$).
  \ii Rank $2$, nullity $2$, $2+2=4=n$.
\end{steps}
\end{prob}

\begin{prob}\textbf{P2.}\ Three vectors in $\RR^3$: by the basis theorem it is enough to check independence. $\det\begin{bmatrix}1&2&1\\0&1&3\\-1&0&4\end{bmatrix}=1(4-0)-2(0+3)+1(0+1)=-1\ne0$. Yes, a basis.
\end{prob}

\begin{prob}\textbf{P3.}\ $\begin{bmatrix}1&2&0\\0&1&-1\\2&5&-1\end{bmatrix}\to\begin{bmatrix}1&0&2\\0&1&-1\\0&0&0\end{bmatrix}$. Pivot columns 1, 2: basis $\{(1,0,2),(2,1,5)\}$, $\dim H=2$. (Indeed $(0,-1,-1)=-2(1,0,2)+(2,1,5)$.)
\end{prob}

\begin{prob}\textbf{P4.}\par\idea{Turn $p(1)=0$ into an equation on the coefficients.}
$p=a_0+a_1t+a_2t^2+a_3t^3$ and $p(1)=a_0+a_1+a_2+a_3=0$, so $a_0=-a_1-a_2-a_3$ and $p=a_1(t-1)+a_2(t^2-1)+a_3(t^3-1)$. So $\{t-1,\ t^2-1,\ t^3-1\}$ spans the subspace. It is independent: the three polynomials have different degrees, so in a relation the highest-degree term with nonzero weight could not cancel. Dimension $3$.
\end{prob}

\begin{prob}\textbf{P5.}\ (a) $\dim\Nul A=8-5=3$. $\Col A$ is a 5-dimensional subspace of $\RR^5$, so it is $\RR^5$ (basis theorem). (b) No: $\rank A\le4$, so $\dim\Nul A=6-\rank A\ge2$.
\end{prob}

\begin{prob}\textbf{P6.}\ Row reduce $[\,\vv v_1\ \vv v_2\ \vv e_1\ \vv e_2\ \vv e_3\,]=\begin{bmatrix}1&0&1&0&0\\2&1&0&1&0\\0&3&0&0&1\end{bmatrix}$: $R_2-2R_1$ gives $(0,1,-2,1,0)$; $R_3-3R_2$ gives $(0,0,6,-3,1)$. Pivot columns 1, 2, 3: basis $\{(1,2,0),(0,1,3),(1,0,0)\}$.
\end{prob}

\begin{prob}\textbf{P7.}
\begin{steps}
  \ii (a) True: more vectors than a basis.
  \ii (b) False: they could be dependent (for example equal), spanning a line or $\{\vv0\}$.
  \ii (c) True: both equal the number of pivots.
  \ii (d) True: $\Nul A\ne\{\vv0\}$, so $A$ is not invertible, so $\det A=0$.
\end{steps}
\end{prob}

\end{document}
